\chapter*{Conclusion}
\addcontentsline{toc}{chapter}{Conclusion}

This thesis presented DBMS-Nautilus, a grammar-based greybox fuzzing system for automated vulnerability detection in SQLite. The central contribution is a grammar engineering methodology based on structural primitives: SQL patterns extracted from CVE root-cause analysis are encoded as composable, weighted production rules without hardcoding any proof-of-concept inputs. The cross-pollination property of this design --- where a single non-terminal serves multiple CVEs simultaneously and the fuzzer freely composes patterns across independent non-terminal groups --- enables genuine compositional discovery of vulnerability-triggering inputs.

The grammar evolved through four versions (v3.0 to v3.3, from 449 to 520 production rules), each refined based on experimental evidence. The key transition occurred at v3.2, which added window function and self-referential generated column primitives, expanding CVE reachability from 2/6 to 6/6. This evolution demonstrated that grammar engineering operates on two orthogonal axes --- weight tuning and structural expansion --- and that both are necessary for effective vulnerability discovery.

Regarding the first research question, DBMS-Nautilus successfully rediscovered three of six target CVEs through compositional grammar generation: CVE-2020-13434 (integer overflow in printf), CVE-2020-9327 (null pointer dereference in generated column resolution), and CVE-2020-13871 (use-after-free in window function processing). None of these were found by embedding proof-of-concept inputs --- the fuzzer independently composed the triggering SQL from separate structural primitives. The three unfound CVEs remained elusive due to combinatorial barriers (CVE-2020-13435 requires five simultaneous structural elements), oracle limitations (CVE-2019-19646 causes an infinite loop rather than a crash), or specific nesting requirements (CVE-2020-15358).

Regarding the second research question, DBMS-Nautilus discovered seven previously unreported bug classes across 89 unique crash instances spanning four SQLite versions. These include FTS5 null pointer dereferences, pervasive float cast overflows, numeric conversion null dereferences, and window function lifecycle errors. The discovery of bugs beyond the grammar's explicit design targets confirms the cross-pollination hypothesis: structural patterns intended for specific CVEs also exercise adjacent code paths containing previously undocumented defects.

Several limitations bound these findings. All experiments target a single DBMS (SQLite), and the results may not generalize to larger server-based systems. The context-free grammar cannot enforce semantic constraints, meaning a fraction of generated inputs are rejected at runtime. The sanitizer-based oracle misses non-crash bugs such as infinite loops and incorrect query results. The grammar design process requires manual domain expertise in both SQL internals and vulnerability patterns.

Four directions warrant further investigation. First, reinforcement learning integration could enable adaptive weight tuning based on coverage and crash-frequency signals, potentially accelerating discovery of bugs requiring rare compositional combinations. Second, cross-DBMS grammar portability should be tested by adapting the structural primitives methodology to MySQL and PostgreSQL. Third, semantic-aware generation extending the grammar with type constraints would reduce wasted executions on semantically invalid inputs. Fourth, automated grammar expansion using large language model analysis of vulnerability disclosures could partially automate the grammar design process.
